% Options for packages loaded elsewhere
\PassOptionsToPackage{unicode}{hyperref}
\PassOptionsToPackage{hyphens}{url}
\documentclass[
]{article}
\usepackage{xcolor}
\usepackage{amsmath,amssymb}
\setcounter{secnumdepth}{-\maxdimen} % remove section numbering
\usepackage{iftex}
\ifPDFTeX
  \usepackage[T1]{fontenc}
  \usepackage[utf8]{inputenc}
  \usepackage{textcomp} % provide euro and other symbols
\else % if luatex or xetex
  \usepackage{unicode-math} % this also loads fontspec
  \defaultfontfeatures{Scale=MatchLowercase}
  \defaultfontfeatures[\rmfamily]{Ligatures=TeX,Scale=1}
\fi
\usepackage{lmodern}
\ifPDFTeX\else
  % xetex/luatex font selection
\fi
% Use upquote if available, for straight quotes in verbatim environments
\IfFileExists{upquote.sty}{\usepackage{upquote}}{}
\IfFileExists{microtype.sty}{% use microtype if available
  \usepackage[]{microtype}
  \UseMicrotypeSet[protrusion]{basicmath} % disable protrusion for tt fonts
}{}
\makeatletter
\@ifundefined{KOMAClassName}{% if non-KOMA class
  \IfFileExists{parskip.sty}{%
    \usepackage{parskip}
  }{% else
    \setlength{\parindent}{0pt}
    \setlength{\parskip}{6pt plus 2pt minus 1pt}}
}{% if KOMA class
  \KOMAoptions{parskip=half}}
\makeatother
\usepackage{color}
\usepackage{fancyvrb}
\newcommand{\VerbBar}{|}
\newcommand{\VERB}{\Verb[commandchars=\\\{\}]}
\DefineVerbatimEnvironment{Highlighting}{Verbatim}{commandchars=\\\{\}}
% Add ',fontsize=\small' for more characters per line
\newenvironment{Shaded}{}{}
\newcommand{\AlertTok}[1]{\textcolor[rgb]{1.00,0.00,0.00}{\textbf{#1}}}
\newcommand{\AnnotationTok}[1]{\textcolor[rgb]{0.38,0.63,0.69}{\textbf{\textit{#1}}}}
\newcommand{\AttributeTok}[1]{\textcolor[rgb]{0.49,0.56,0.16}{#1}}
\newcommand{\BaseNTok}[1]{\textcolor[rgb]{0.25,0.63,0.44}{#1}}
\newcommand{\BuiltInTok}[1]{\textcolor[rgb]{0.00,0.50,0.00}{#1}}
\newcommand{\CharTok}[1]{\textcolor[rgb]{0.25,0.44,0.63}{#1}}
\newcommand{\CommentTok}[1]{\textcolor[rgb]{0.38,0.63,0.69}{\textit{#1}}}
\newcommand{\CommentVarTok}[1]{\textcolor[rgb]{0.38,0.63,0.69}{\textbf{\textit{#1}}}}
\newcommand{\ConstantTok}[1]{\textcolor[rgb]{0.53,0.00,0.00}{#1}}
\newcommand{\ControlFlowTok}[1]{\textcolor[rgb]{0.00,0.44,0.13}{\textbf{#1}}}
\newcommand{\DataTypeTok}[1]{\textcolor[rgb]{0.56,0.13,0.00}{#1}}
\newcommand{\DecValTok}[1]{\textcolor[rgb]{0.25,0.63,0.44}{#1}}
\newcommand{\DocumentationTok}[1]{\textcolor[rgb]{0.73,0.13,0.13}{\textit{#1}}}
\newcommand{\ErrorTok}[1]{\textcolor[rgb]{1.00,0.00,0.00}{\textbf{#1}}}
\newcommand{\ExtensionTok}[1]{#1}
\newcommand{\FloatTok}[1]{\textcolor[rgb]{0.25,0.63,0.44}{#1}}
\newcommand{\FunctionTok}[1]{\textcolor[rgb]{0.02,0.16,0.49}{#1}}
\newcommand{\ImportTok}[1]{\textcolor[rgb]{0.00,0.50,0.00}{\textbf{#1}}}
\newcommand{\InformationTok}[1]{\textcolor[rgb]{0.38,0.63,0.69}{\textbf{\textit{#1}}}}
\newcommand{\KeywordTok}[1]{\textcolor[rgb]{0.00,0.44,0.13}{\textbf{#1}}}
\newcommand{\NormalTok}[1]{#1}
\newcommand{\OperatorTok}[1]{\textcolor[rgb]{0.40,0.40,0.40}{#1}}
\newcommand{\OtherTok}[1]{\textcolor[rgb]{0.00,0.44,0.13}{#1}}
\newcommand{\PreprocessorTok}[1]{\textcolor[rgb]{0.74,0.48,0.00}{#1}}
\newcommand{\RegionMarkerTok}[1]{#1}
\newcommand{\SpecialCharTok}[1]{\textcolor[rgb]{0.25,0.44,0.63}{#1}}
\newcommand{\SpecialStringTok}[1]{\textcolor[rgb]{0.73,0.40,0.53}{#1}}
\newcommand{\StringTok}[1]{\textcolor[rgb]{0.25,0.44,0.63}{#1}}
\newcommand{\VariableTok}[1]{\textcolor[rgb]{0.10,0.09,0.49}{#1}}
\newcommand{\VerbatimStringTok}[1]{\textcolor[rgb]{0.25,0.44,0.63}{#1}}
\newcommand{\WarningTok}[1]{\textcolor[rgb]{0.38,0.63,0.69}{\textbf{\textit{#1}}}}
\usepackage{longtable,booktabs,array}
\newcounter{none} % for unnumbered tables
\usepackage{calc} % for calculating minipage widths
% Correct order of tables after \paragraph or \subparagraph
\usepackage{etoolbox}
\makeatletter
\patchcmd\longtable{\par}{\if@noskipsec\mbox{}\fi\par}{}{}
\makeatother
% Allow footnotes in longtable head/foot
\IfFileExists{footnotehyper.sty}{\usepackage{footnotehyper}}{\usepackage{footnote}}
\makesavenoteenv{longtable}
\setlength{\emergencystretch}{3em} % prevent overfull lines
\providecommand{\tightlist}{%
  \setlength{\itemsep}{0pt}\setlength{\parskip}{0pt}}
\usepackage{bookmark}
\IfFileExists{xurl.sty}{\usepackage{xurl}}{} % add URL line breaks if available
\urlstyle{same}
\hypersetup{
  pdftitle={AI World: A Scenario-Pack Framework for Verifiable Multi-Agent Problem Solving},
  pdfauthor={Bingbing Zhang},
  hidelinks,
  pdfcreator={LaTeX via pandoc}}

\title{AI World: A Scenario-Pack Framework for Verifiable Multi-Agent
Problem Solving}
\author{Bingbing Zhang\\Independent Researcher, TraceArena Open-source Project\\\texttt{tonyhyworld@gmail.com}}
\date{2026}

\begin{document}
\maketitle

\section{AI World: A Scenario-Pack Framework for Verifiable Multi-Agent
Problem
Solving}\label{ai-world-a-scenario-pack-framework-for-verifiable-multi-agent-problem-solving}

\subsection{Abstract}\label{abstract}

Large language models are increasingly used as agents, yet most agent
systems are built around one-shot answers or fixed workflows. This makes
it difficult for domain experts to encode their own constraints, for
multiple agents to explore competing strategies, and for users to verify
how an outcome was produced. We present \textbf{AI World}, implemented
by the open-source \textbf{TraceArena} runtime, a scenario-pack
framework for turning domain problems into executable worlds. A scenario
pack declares goals, roles, resources, rules, tools, external data,
termination conditions, settlement logic, and presentation metadata.
Within a loaded world, multiple agents execute a bounded
observation-\/-planning-\/-action-\/-feedback loop. A world kernel
validates actions and advances state; a settlement runtime computes
outcomes from simulation rules, external facts, deterministic verifiers,
or hybrid combinations; and an event ledger preserves an auditable trace
that can be replayed visually. The contribution is a reusable interface
between domain expertise and agent execution: experts define the world,
while the runtime provides orchestration, verification, and
observability. We describe the architecture, execution protocol,
scenario-pack design, and an evaluation methodology for comparing agents
and strategies under identical constraints. The system is released with
public examples and reproducibility guidance.

\textbf{Keywords:} multi-agent systems, agent environments, scenario
packs, agent loop, verifiable execution, AI simulation, decision
traceability, visual replay

\subsection{1. Introduction}\label{1-introduction}

An agent that produces a plausible paragraph is not necessarily an agent
that can solve a real problem. Real decisions unfold over time, consume
resources, face permissions and constraints, interact with other
decision makers, and produce consequences that may only become visible
later. Domain experts also need to express what counts as a legal action
and what counts as a good result. These requirements are difficult to
satisfy with a single prompt, a one-shot benchmark, or a fixed workflow
engine.

AI World addresses this gap by providing an executable world abstraction
for domain experts. A world is defined by a scenario pack rather than
hard-coded into the operating system. The expert declares the domain
objective, participants, resources, rules, tools, evidence sources, and
settlement criteria. Multiple agents can then act in the same world,
compete or cooperate, receive feedback, revise plans, and continue until
a goal or termination condition is reached. The runtime makes the entire
chain visible: what an agent observed, what it decided, what action it
attempted, how the world responded, and how the final result was
settled.

Our central thesis is:

\begin{quote}
A domain expert should be able to turn a real problem into a runnable
world without first building a complete agent platform, while every
important action remains verifiable and replayable.
\end{quote}

This paper makes four contributions:

\begin{enumerate}
\def\labelenumi{\arabic{enumi}.}
\tightlist
\item
  A scenario-pack abstraction that separates domain definitions from the
  general-purpose runtime.
\item
  A multi-agent execution loop with explicit world validation, feedback,
  bounded retries, and termination.
\item
  A settlement and trace model that supports simulated,
  external-reality, deterministic-verifier, and hybrid worlds.
\item
  A visual presentation path that converts verified facts into
  human-readable, replayable decision stories without making the
  narrator authoritative over the world state.
\end{enumerate}

\subsection{2. Design Goals and
Non-Goals}\label{2-design-goals-and-non-goals}

\subsubsection{2.1 Goals}\label{21-goals}

AI World is designed to:

\begin{itemize}
\tightlist
\item
  lower the barrier for domain experts to build agent environments;
\item
  support continuous action rather than one-shot completion;
\item
  allow multiple agents and strategies to be compared under the same
  world constraints;
\item
  keep world state, actions, events, settlement, and traces auditable;
\item
  work with different LLM providers and internal or external tools;
\item
  make complex execution understandable through live visualization and
  replay.
\end{itemize}

\subsubsection{2.2 Non-goals}\label{22-non-goals}

AI World is not a claim that an LLM is correct, safe, or autonomous in
every setting. It is not a replacement for production controls,
regulatory review, human authorization, or physical safety systems. It
is also not primarily a dataset marketplace. Trajectories and evaluation
records are useful outputs, but the product\textquotesingle s primary
value is providing a reusable world framework for solving and inspecting
domain problems.

\subsection{3. System Model}\label{3-system-model}

Let a scenario pack define a world (W) with state (s\_t), agents (A),
actions (U), tools (T), observations (O), and a settlement function (S).
At step (t), agent (a\_i) receives an authorized observation (o\_\{i,t\}
= V\_i(s\_t)), where (V\_i) is the visibility policy for that role. The
agent produces a plan and a candidate action (u\_\{i,t\}). The world
kernel applies a legality function (L(u\_\{i,t\}, s\_t)). If the action
is legal, the transition function produces (s\_\{t+1\}) and a world
event (e\_t); otherwise the kernel returns a structured failure that can
be used by the agent for reflection.

The run ends when a goal predicate, a termination predicate, or a
resource/time budget is reached. The settlement runtime computes (r =
S(s\_\{0:T\}, e\_\{0:T\}, x)), where (x) may include external facts such
as market prices or deterministic verifier outputs. The trace ledger
stores observations, model decisions, tool calls, actions, failures,
events, settlement records, and presentation facts.

\subsection{4. Architecture}\label{4-architecture}

\begin{Shaded}
\begin{Highlighting}[]
\NormalTok{Scenario Pack}
\NormalTok{  goals / roles / resources / rules / tools / settlement / presentation}
\NormalTok{                              |}
\NormalTok{                              v}
\NormalTok{TraceArena Runtime}
\NormalTok{  loader {-}\textgreater{} scheduler {-}\textgreater{} agent harness {-}\textgreater{} world kernel {-}\textgreater{} event ledger}
\NormalTok{                                      |                 |}
\NormalTok{                                      v                 v}
\NormalTok{                               LLM providers       settlement runtime}
\NormalTok{                                      \textbackslash{}                 /}
\NormalTok{                                       v               v}
\NormalTok{                                  verified facts and traces}
\NormalTok{                                             |}
\NormalTok{                                             v}
\NormalTok{                              director {-}\textgreater{} renderer {-}\textgreater{} replay / console}
\end{Highlighting}
\end{Shaded}

\subsubsection{4.1 Scenario-pack layer}\label{41-scenario-pack-layer}

The scenario pack is the extension boundary. It declares the world
vocabulary, roles, actions, resources, tools, clocks, goals, settlement
manifest, and rendering bindings. The runtime does not assume that every
world has the same action schema or reward function.

\subsubsection{4.2 Agent harness}\label{42-agent-harness}

The harness converts an LLM into a bounded actor. It assembles role
prompts and world observations, exposes approved capabilities, invokes
tools, parses structured actions, collects feedback, and controls
retries and budgets. A failed action is not silently treated as success;
it becomes an explicit feedback event.

\subsubsection{4.3 World kernel}\label{43-world-kernel}

The world kernel is authoritative for state transitions. It enforces
permissions, action schemas, resource availability, visibility, phase
rules, and event creation. This separation prevents a language model or
a narrative component from directly declaring its own outcome.

\subsubsection{4.4 Settlement runtime}\label{44-settlement-runtime}

Different domains require different sources of truth. AI World supports
four settlement modes:

{\def\LTcaptype{none} % do not increment counter
\begin{longtable}[]{@{}lll@{}}
\toprule\noalign{}
Mode & Source of outcome & Example \\
\midrule\noalign{}
\endhead
\bottomrule\noalign{}
\endlastfoot
Simulation & Scenario rules & social influence, resource changes \\
External reality & External data or service & observed prices, external
task result \\
Deterministic verifier & Reproducible checker & tests, schemas, exact
graders \\
Hybrid & External facts plus deterministic rules & market price plus
portfolio ledger \\
\end{longtable}
}

\subsubsection{4.5 Trace ledger and
replay}\label{45-trace-ledger-and-replay}

The ledger records the facts needed to explain a run. A run can be
inspected as a timeline or replayed from fixtures. This supports
debugging, model comparison, expert review, and later experimental
analysis.

\subsubsection{4.6 Director and
renderer}\label{46-director-and-renderer}

The director translates verified events and settlement records into
audience-facing descriptions and rendering commands. It may explain why
an action mattered, but it cannot replace the world kernel or rewrite
the settlement record. This design separates narrative clarity from
factual authority.

\subsection{5. Agent Loop Protocol}\label{5-agent-loop-protocol}

The canonical loop is:

\begin{Shaded}
\begin{Highlighting}[]
\NormalTok{observe {-}\textgreater{} plan {-}\textgreater{} select capability {-}\textgreater{} execute {-}\textgreater{} validate}
\NormalTok{    \^{}                                         |}
\NormalTok{    |{-}{-}{-}{-}{-}{-}{-}{-}{-}{-}{-}{-}{-} feedback / reflect {-}{-}{-}{-}{-}{-}{-}{-}|}
\end{Highlighting}
\end{Shaded}

The loop is bounded by explicit policies:

\begin{itemize}
\tightlist
\item
  maximum steps or ticks;
\item
  token and tool budgets;
\item
  action permissions;
\item
  world termination conditions;
\item
  retry and failure policies;
\item
  human approval gates where required.
\end{itemize}

This allows ``continuous execution'' to mean sustained goal-directed
work rather than an unbounded or uncontrolled process.

\subsection{6. Building a Domain Scenario
Pack}\label{6-building-a-domain-scenario-pack}

A domain team can create a scenario pack in five stages:

\begin{enumerate}
\def\labelenumi{\arabic{enumi}.}
\tightlist
\item
  \textbf{State the problem.} Define a measurable objective and what
  constitutes completion.
\item
  \textbf{Model the world.} Identify actors, resources, time,
  information, constraints, and external dependencies.
\item
  \textbf{Declare actions and tools.} Specify which actions are legal,
  which tools each role can use, and which side effects require
  authorization.
\item
  \textbf{Define settlement.} Select simulation rules, an external fact
  source, a deterministic verifier, or a hybrid.
\item
  \textbf{Define presentation.} Choose the events, metrics, decisions,
  and comparisons that experts and viewers need to see.
\end{enumerate}

The resulting pack can be run repeatedly with different models, prompts,
seeds, role policies, or tool configurations while keeping the world
definition fixed.

\subsection{7. Applications}\label{7-applications}

AI World is domain-neutral. Public examples and target applications
include:

\begin{itemize}
\tightlist
\item
  \textbf{Capital markets:} compare research and portfolio strategies
  under market data, risk budgets, and trading rules.
\item
  \textbf{City governance:} test transport, emergency response, and
  resource allocation policies with multiple stakeholders.
\item
  \textbf{Drug discovery:} plan experiments under budget, evidence, and
  trial-stage constraints.
\item
  \textbf{Logistics:} coordinate orders, inventory, routes, weather, and
  disruptions.
\item
  \textbf{Manufacturing and robotics:} schedule production, handle
  faults, and respect safety boundaries in a controlled world.
\item
  \textbf{Education and research:} create reproducible experimental
  worlds for teaching, hypothesis testing, and model evaluation.
\end{itemize}

\subsection{8. Evaluation Methodology}\label{8-evaluation-methodology}

This paper intentionally distinguishes implemented capabilities from
measured results. A reproducible evaluation should include:

\subsubsection{8.1 Research questions}\label{81-research-questions}

\begin{itemize}
\tightlist
\item
  Does explicit world validation reduce illegal or impossible actions?
\item
  Does multi-agent competition discover different and better paths than
  a single agent?
\item
  Does a feedback-and-reflection loop improve recovery after failure?
\item
  Does visual replay improve expert ability to identify why a result
  occurred?
\item
  How stable are results across models, seeds, and repeated runs?
\end{itemize}

\subsubsection{8.2 Conditions}\label{82-conditions}

Compare at least:

\begin{enumerate}
\def\labelenumi{\arabic{enumi}.}
\tightlist
\item
  single agent vs. multiple agents;
\item
  one-shot response vs. bounded Agent Loop;
\item
  no tools vs. approved tools;
\item
  no reflection vs. feedback-triggered reflection;
\item
  narrative-only judgment vs. authoritative settlement.
\end{enumerate}

\subsubsection{8.3 Metrics}\label{83-metrics}

Recommended metrics include goal completion, settlement score, resource
cost, rule-violation rate, recovery rate after failed actions, number of
tool calls, run-to-run variance, trace completeness, and expert review
time.

All reported numbers should include the scenario-pack version, runtime
commit, model and provider, prompts, tool configuration, budget, seed,
hardware, and raw replay artifacts.

\subsection{9. Limitations and Risks}\label{9-limitations-and-risks}

The quality of a world depends on the quality of its rules, data, tools,
and settlement design. A poorly specified scenario can produce
precise-looking but invalid conclusions. External data can be delayed,
unavailable, or subject to licensing constraints. LLM behavior remains
probabilistic and can be sensitive to prompts and provider changes. For
high-impact domains, AI World should be used for analysis, simulation,
and decision support with appropriate human oversight, not as an
unreviewed authority.

\subsection{10. Open-source Release}\label{10-open-source-release}

TraceArena is released as the open-source runtime for AI World. The
repository includes a scenario-pack template, runnable examples, local
replay paths, deployment documentation, and contribution guidance. The
community can extend the system by contributing domain packs rather than
modifying the core runtime for every new use case.

\subsection{11. Conclusion}\label{11-conclusion}

AI World provides a practical interface between domain expertise and
multi-agent execution. It allows a professional team to define a world,
load it into a general runtime, let agents act and compete under
explicit constraints, and inspect the path to the outcome. The central
object is not a conversational answer but a verifiable run: a sequence
of observations, decisions, actions, feedback, events, and settlement
records that can be watched, replayed, compared, and improved.

The long-term vision is an ecosystem of domain-specific AI Worlds built
on a shared open-source runtime. In this model, a financial expert, a
city planner, a medical researcher, or a logistics operator can
contribute a world in which AI works on the problems that matter to that
field.

\subsection{References}\label{references}

{[}1{]} Yao et al. ReAct: Synergizing Reasoning and Acting in Language
Models. 2023.\\
{[}2{]} Wu et al. AutoGen: Enabling Next-Gen LLM Applications. 2023.\\
{[}3{]} Wang et al. Voyager: An Open-Ended Embodied Agent with Large
Language Models. 2023.\\
{[}4{]} Zhou et al. WebArena: A Realistic Web Environment for Building
Autonomous Agents. 2024.\\
{[}5{]} TraceArena contributors. TraceArena: Open-source AI World
runtime and scenario packs. 2026.
\url{https://github.com/tonyhyworld/TraceArena}

\end{document}
